% ═══════════════════════════════════════════════════════════════
% LERNBLATT - Physik Klasse 6: Bewegungsarten erkennen
% Stunde 1 von 4 | Doppelseiten-Format
% ═══════════════════════════════════════════════════════════════

\documentclass[11pt, a4paper]{article}

\usepackage[utf8]{inputenc}
\usepackage[T1]{fontenc}
\usepackage[ngerman]{babel}
\usepackage{helvet}
\renewcommand{\familydefault}{\sfdefault}

\usepackage[margin=1.5cm, top=1.8cm, bottom=1.5cm]{geometry}
\usepackage{enumitem}
\usepackage{tabularx}
\usepackage{array}
\usepackage{multicol}
\usepackage{tikz}
\usepackage{amssymb}
\usepackage{fancyhdr}
\usepackage{lastpage}
\usepackage{xcolor}
\usepackage{tcolorbox}
\tcbuselibrary{skins, breakable}

% Farben
\definecolor{physik}{HTML}{2c5aa0}
\definecolor{physik-hell}{HTML}{e8f0fa}
\definecolor{hellgrau}{HTML}{F5F5F5}
\definecolor{rahmen}{HTML}{CCCCCC}

% Kopfzeile
\pagestyle{fancy}
\fancyhf{}
\fancyhead[L]{\small\textcolor{physik}{\textbf{Physik}} | Bewegungsarten}
\fancyhead[R]{\small Klasse 6}
\fancyfoot[C]{\small\thepage}
\renewcommand{\headrulewidth}{0pt}

% Boxen
\newtcolorbox{merkkasten}[1][Merke]{
    colback=physik-hell,
    colframe=physik,
    boxrule=1.5pt,
    arc=0pt,
    left=8pt, right=8pt, top=8pt, bottom=6pt,
    title={\small\textbf{#1}},
    coltitle=white,
    fonttitle=\sffamily,
    attach boxed title to top left={yshift=-2mm, xshift=5mm},
    boxed title style={colback=physik, arc=0pt}
}

\newtcolorbox{infobox}{
    colback=hellgrau,
    colframe=hellgrau,
    boxrule=0pt,
    arc=0pt,
    left=8pt, right=8pt, top=6pt, bottom=6pt,
    borderline west={3pt}{0pt}{physik}
}

% Befehle
\newcommand{\seitentitel}[1]{%
    \noindent\textcolor{physik}{\rule{\textwidth}{2pt}}\\[0.3em]
    {\LARGE\bfseries\textcolor{physik}{#1}}\\[0.2em]
    \textcolor{physik}{\rule{\textwidth}{0.5pt}}
    \vspace{0.5em}
}

\newcommand{\abschnitt}[1]{%
    \vspace{8pt}%
    \noindent{\large\bfseries\textcolor{physik}{#1}}%
    \vspace{4pt}%
}

\newcommand{\luecke}[1]{\underline{\hspace{#1}}}

\setlength{\parindent}{0pt}
\setlength{\parskip}{4pt}

\begin{document}

% ═══════════════════════════════════════════════════════════════
% SEITE 1: THEORIE
% ═══════════════════════════════════════════════════════════════

\seitentitel{Bewegungsarten erkennen}

\begin{infobox}
\textbf{Einstieg:} Schau dich um – überall gibt es Bewegungen! Ein Auto fährt, ein Karussell dreht sich, eine Schlange kriecht. Aber wie unterscheiden Physiker diese Bewegungen?
\end{infobox}

\abschnitt{1. Was ist Bewegung?}

\begin{merkkasten}[Kasten 1: Bewegung]
Ein Körper ist in \textbf{Bewegung}, wenn er seinen \luecke{1.8cm} verändert.

Wir beschreiben Bewegungen mit zwei Eigenschaften:
\begin{itemize}[leftmargin=1.5em, itemsep=1pt]
    \item Die \textbf{Bahnform} beschreibt, wie der \luecke{1.8cm} aussieht.
    \item Die \textbf{Bewegungsart} beschreibt, ob sich die \luecke{2.5cm} ändert.
\end{itemize}
\end{merkkasten}

\abschnitt{2. Einteilung nach Bahnform}

\begin{center}
\begin{tikzpicture}[scale=0.8]
    % Geradlinig
    \draw[thick, physik, ->] (0,0) -- (3,0);
    \node[below] at (1.5,-0.3) {\small\textbf{geradlinig}};
    \node[above, font=\scriptsize] at (1.5,0.3) {Auto auf Autobahn};

    % Krummlinig (Oberbegriff)
    \node[below] at (8,-0.3) {\small\textbf{krummlinig}};

    % Kreisförmig (Sonderfall)
    \draw[thick, physik, ->] (5.5,0) arc (180:-180:0.8);
    \node[above, font=\scriptsize] at (6.3,0.8) {Karussell};
    \node[below, font=\scriptsize] at (6.3,-1.1) {(kreisförmig)};

    % Beliebig krummlinig
    \draw[thick, physik, ->] (8.5,0) .. controls (9.5,0.8) and (10.5,-0.5) .. (11.5,0.3);
    \node[above, font=\scriptsize] at (10,0.7) {Schlange};
    \node[below, font=\scriptsize] at (10,-0.6) {(beliebig)};
\end{tikzpicture}
\end{center}

\begin{merkkasten}[Kasten 2: Bahnformen]
Es gibt zwei Hauptformen:
\begin{itemize}[leftmargin=1.5em, itemsep=2pt]
    \item \textbf{Geradlinig:} Der Körper bewegt sich auf einer \luecke{2.5cm}.
    \item \textbf{Krummlinig:} Der Körper bewegt sich auf einer \luecke{2.5cm} Bahn.
\end{itemize}

\textit{Sonderfall:} Die \textbf{kreisförmige} Bewegung ist eine spezielle krummlinige Bewegung -- der Körper bewegt sich auf einem \luecke{2cm}.
\end{merkkasten}

\abschnitt{3. Einteilung nach Bewegungsart}

\begin{center}
\begin{tikzpicture}[scale=0.8]
    % Gleichförmig
    \foreach \x in {0,1,2,3} {
        \fill[physik] (\x,0) circle (4pt);
    }
    \draw[->] (0,-0.5) -- (3.5,-0.5);
    \node[below] at (1.75,-0.7) {\small\textbf{gleichförmig}};
    \node[above, text width=3cm, align=center, font=\scriptsize] at (1.5,0.4) {Gleiche Abstände\\= gleiche Geschwindigkeit};

    % Ungleichförmig
    \foreach \x/\pos in {0/6, 0.3/6.3, 0.8/6.8, 1.5/7.5, 2.5/8.5} {
        \fill[physik] (\pos,0) circle (4pt);
    }
    \draw[->] (6,-0.5) -- (9,-0.5);
    \node[below] at (7.5,-0.7) {\small\textbf{ungleichförmig}};
    \node[above, text width=3cm, align=center, font=\scriptsize] at (7.5,0.4) {Verschiedene Abstände\\= Geschwindigkeit ändert sich};
\end{tikzpicture}
\end{center}

\begin{merkkasten}[Kasten 3: Bewegungsarten]
Es gibt zwei Bewegungsarten:
\begin{itemize}[leftmargin=1.5em, itemsep=2pt]
    \item \textbf{Gleichförmig:} Die Geschwindigkeit bleibt \luecke{2cm}.
    \item \textbf{Ungleichförmig:} Die Geschwindigkeit \luecke{2cm} sich.
\end{itemize}
\textit{Beispiele für ungleichförmig:} Anfahren, Bremsen, Beschleunigen
\end{merkkasten}

% ═══════════════════════════════════════════════════════════════
% SEITE 2: ANWENDUNG
% ═══════════════════════════════════════════════════════════════

\newpage

\seitentitel{Bewegungsarten erkennen}

\abschnitt{4. Beispiele aus dem Alltag}

\begin{center}
\renewcommand{\arraystretch}{1.4}
\begin{tabular}{|p{4.5cm}|c|c|}
\hline
\textbf{Beispiel} & \textbf{Bahnform} & \textbf{Bewegungsart} \\
\hline
Zug auf gerader Strecke & geradlinig & gleichförmig \\
\hline
Fahrrad beim Anfahren & geradlinig & ungleichförmig \\
\hline
Gondel am Riesenrad & krummlinig (kreisförmig) & gleichförmig \\
\hline
Achterbahn in der Kurve & krummlinig & ungleichförmig \\
\hline
Sekundenzeiger einer Uhr & krummlinig (kreisförmig) & gleichförmig \\
\hline
\end{tabular}
\end{center}

\vspace{0.5em}

\abschnitt{5. Übung: Ordne zu!}

Welche Bahnform und Bewegungsart haben diese Bewegungen? Kreuze an.

\textit{Hinweis: Kreisförmige Bewegungen sind krummlinig!}

\begin{center}
\renewcommand{\arraystretch}{1.6}
\begin{tabular}{|p{3.5cm}|c|c|c|c|}
\hline
& \multicolumn{2}{c|}{\textbf{Bahnform}} & \multicolumn{2}{c|}{\textbf{Bewegungsart}} \\
\cline{2-5}
\textbf{Beispiel} & gerad- & krumm- & gleich- & ungleich- \\
& linig & linig & förmig & förmig \\
\hline
a) Rolltreppe & $\square$ & $\square$ & $\square$ & $\square$ \\
\hline
b) Schaukel & $\square$ & $\square$ & $\square$ & $\square$ \\
\hline
c) Aufzug (volle Fahrt) & $\square$ & $\square$ & $\square$ & $\square$ \\
\hline
d) Fußball beim Schuss & $\square$ & $\square$ & $\square$ & $\square$ \\
\hline
\end{tabular}
\end{center}

\vspace{0.5em}

\abschnitt{6. Zusammenfassung}

\begin{merkkasten}[Das Wichtigste]
\textbf{Bewegung} = Ein Körper verändert seinen Ort.

\textbf{Bahnform} = Wie sieht der Weg aus?
\begin{itemize}[leftmargin=1.5em, itemsep=1pt]
    \item Geradlinig (gerade Linie)
    \item Krummlinig (gebogene Bahn) -- dazu gehört auch kreisförmig!
\end{itemize}

\textbf{Bewegungsart} = Ändert sich die Geschwindigkeit?
\begin{itemize}[leftmargin=1.5em, itemsep=1pt]
    \item Gleichförmig (bleibt gleich) -- Ungleichförmig (ändert sich)
\end{itemize}
\end{merkkasten}

\end{document}
